\section{Agentic 模型与传统后训练范式的分野}

\subsection{为什么传统 RLHF 不能直接覆盖 Agent 需求}
传统 RLHF 的核心是偏好建模：收集人类比较数据，训练 reward model，再用 SFT、PPO/GRPO 等算法提升 ``被人偏好'' 的概率。这条链路在对话系统上很有效，但在 Agent 场景里会暴露两个缺口：第一，许多任务存在硬约束（例如 JSON 结构、单元测试通过、SQL 正确执行）；第二，任务执行过程会跨越多轮工具调用，最终结果取决于过程正确性而非表面表达。

\begin{warningbox}{常见误区：把 Agent 训练当成 ``更长的对话 RLHF''}
如果只把 Agent 当作更长上下文的聊天任务，会出现三类问题：
\begin{itemize}
    \item 过程错误被最终话术掩盖，导致 ``可读不可用''；
    \item 奖励信号过度依赖主观偏好，难以覆盖硬性正确性；
    \item 训练后模型在真实工具链中不稳定，出现格式漂移和动作失配。
\end{itemize}
\end{warningbox}

\subsection{双目标训练：Preference + Verifiability}
讲者的建议不是抛弃偏好目标，而是建立双目标后训练框架：模型既要保留自然交互能力，又要在环境中完成可验证任务。具体做法通常是让数据集和评估集同时包含 ``偏好维度'' 与 ``任务维度''，并在训练阶段分配不同权重。

\begin{table}[H]
\centering
\begin{tabular}{p{2.8cm}p{4.1cm}p{4.6cm}}
\toprule
\textbf{维度} & \textbf{传统对话后训练} & \textbf{可验证 Agent 后训练} \\
\midrule
训练对象 & 单轮或短多轮文本输出 & 长轨迹动作序列（含工具调用） \\
奖励来源 & 人类偏好模型为主 & verifier + 偏好联合 \\
失败表现 & 不够礼貌或不够有用 & 任务失败、格式失败、状态污染 \\
验证成本 & 高度人工主观判断 & 可程序化自动验证比例更高 \\
\bottomrule
\end{tabular}
\caption{两种后训练范式的核心差异}
\end{table}

\subsection{从 ``回答'' 到 ``行动轨迹''}
课程中的一个关键表述是：Agent 的核心单位不是 token，而是 trajectory。一次任务里模型需要先理解目标，再决定何时调用工具、如何解释工具反馈、是否需要重试、何时终止并汇报。这个过程本质是序列决策问题。

\begin{importantbox}{后训练单位变化}
当训练单位从 ``response'' 变为 ``trajectory'' 后，数据标注、奖励计算、评估指标、甚至日志系统都要改变。很多工程团队的瓶颈并不在模型本身，而在没有把这套基础设施同步升级。
\end{importantbox}

\subsection{本章小结}
本章结论是：Agent 后训练必须从 ``偏好对齐'' 升级为 ``偏好 + 可验证任务完成'' 的联合优化，并把训练单位重定义为可执行轨迹，而非单条文本回复。

